% ============================================================================
% Formatting requirements for future work
% ============================================================================
% 1. Follow the standard template's layout and styling.
% 2. Add every new requirement to this summary in both note projects and the standard template.
% 3. Transcribe notes almost word for word, preserving notation and order; omit crossed-out content.
% 4. Store each section in sections/ and include its .tex file with \input.
% 5. Match source boxes with framed. Use \caseheading only when BOTH conditions are met: the text is obviously a TITLE, and it heads an IMPORTANT group of contents. Ordinary introductory phrases (e.g. The Pauli matrices are) must remain plain text inside the frame; a frame alone does not justify a heading.
% 6. Use \boldsymbol for vector symbols. amsmath provides it; do not add the bm dependency.
% 7. Load mathtools and use dcases instead of cases for display-style case systems.
% 8. Number displayed equations by section. Use equation for a single display; combine consecutive equation blocks with no intervening prose into align. Use align instead of aligned inside equation; suppress intermediate row numbers with \notag when one number is appropriate, leaving the final row numbered. Keep separate numbers for distinct results where useful.
% 9. Use \texorpdfstring for mathematical titles/headings, with a plain-text PDF alternative.
% 10. Indent nested source with four spaces; preserve environment boundaries and matrix/equation rows, keeping optional row spacing attached to \\.
% 11. Compile to stable numbering, contents, and bookmarks with zero warnings; inspect PDF spacing, frames, and clipping.
% 12. Add readable source spaces around operators and &, after commas, and between indexed factors/index commands; keep scripts and arguments attached to their symbols.
% 13. Rely on VS Code word wrap; do not manually wrap prose, comments, or logical equation rows. Keep paragraph breaks and structural line breaks.
% 14. Use \bra{}, \ket{}, and \braket{} for bras, kets, and inner products/matrix elements instead of manually assembling angle delimiters and bars.
% 15. Use consistent title case for section and subsection headings; use plural topic names where appropriate (e.g. One-Particle States and Covariant Fields), and spell out representation in headings.
% 16. Correct minor grammar and unnatural wording while preserving mathematical meaning, notation, and the source order.
% 18. In align, place an implication arrow before the alignment marker: \Rightarrow \quad & ..., so the expression to the right of the arrow aligns with the other rows.
% 19. Use \tensor{R}{_\mu^{\mu'}_\nu^{\nu'}} for ordered tensor indices instead of chaining scripts with empty {} groups; apply the same convention to other tensor symbols.
% 20. Use \frac{}{} for fractions; do not use \dfrac.
% 21. Use ordinary \\ row breaks, not starred \\* row breaks; preserve optional row spacing.
% 22. Read mathematical expressions for meaning when spacing source: separate product factors, operator actions, individual tensor indices, and spin-state labels (e.g. i (* F), \ket{A a; B b}, and s_j s_k X). Fully brace both fraction arguments, e.g. \frac{1}{4} (I + s_j L + s_k R + s_j s_k X); do not use shorthand such as \frac14. Preserve commands, scripts, and mathematical meaning.
% 23. Allow align environments to break naturally across pages; let LaTeX handle pagination. Do not add samepage wrappers or other page-break restrictions to keep an alignment, its heading, or introductory text together.
% ============================================================================

\RequirePackage{fix-cm} % Scalable Computer Modern math at the template heading sizes.
\documentclass[a4paper,12pt]{article}

\usepackage{amsfonts}
\usepackage{amsmath}
\usepackage{mathtools}
\usepackage{amssymb}
\usepackage{array}
\usepackage{framed}
\usepackage[medium]{titlesec}
\usepackage{cite}
\usepackage[normalem]{ulem}
\usepackage{extarrows}
\usepackage{slashed}
\usepackage{isodateo}
\usepackage{graphicx}
\usepackage{capt-of}
\usepackage[dvipsnames]{xcolor}
\usepackage[bookmarksnumbered=true,bookmarksopen=true]{hyperref}
\hypersetup{
    colorlinks,
    linkcolor=NavyBlue,
    citecolor=PineGreen,
    urlcolor=PineGreen
}
\usepackage[hmargin=.7in,vmargin=1.1in]{geometry}
% \usepackage{indentfirst}
\usepackage{booktabs}
\usepackage{bbm}

\linespread{1.1}
% \graphicspath{{fig/}}

% \newcommand{\FR}[2]{\displaystyle\frac{\,{#1}\,}{#2}}
% \newcommand{\fr}[2]{\mbox{$\frac{\,{#1}\,}{#2}$}}
% \newcommand{\n}{\nonumber}
% \renewcommand{\rm}{\mathrm}
\renewcommand{\thefootnote}{\arabic{footnote}}

% \newcommand{\red}{\color{red}}
% \newcommand{\blue}{\color{blue}}
% \newcommand{\gray}{\color{gray}}
% \newcommand{\hl}{\color[rgb]{1,0,.5}}

% \def\bge{\begin{equation}}
% \def \ede{\end{equation}}
% \def\bga{\begin{aligned}}
% \def\eda{\end{aligned}}
% \def\bgb{\begin{bmatrix}}
% \def\edb{\end{bmatrix}}
% \def\bgp{\begin{pmatrix}}
% \def\edp{\end{pmatrix}}
% \def\bgm{\begin{matrix}}
% \def\edm{\end{matrix}}
% \def\bgs{\begin{subequations}}
% \def\eds{\end{subequations}}

% \newcommand{\order}[1]{\mathcal{O}({#1})}
% \def\di{{\mathrm{d}}}
% \def\D{{\mathrm{D}}}
% \def\Di{{\mathcal{D}}}

% \def\mb{\mathbf}
% \def\ms{\mathscr}
% \def\mf{\mathfrak}
% \def\mc{\mathcal}
% \def\cl{\mathrm{cl}}
% \def\pd{\partial}
% \def\ld{{\mathscr{L}}}
% \def\hd{{\mathscr{H}}}
% \def\z{{\bar{z}}}
% \def\sla{\slashed}
% \def\const{\mathrm{const.}}
% \setlength\unitlength{1mm}
% \def\tr{\mathrm{\,tr\,}}
% \def\Tr{\mathrm{\,Tr\,}}
% \def\Det{\mathrm{\,Det\,}}
% \def\to{\rightarrow}
% \def\To{\Rightarrow}
% \def\ii{\mathrm{i}}

% \def\al{\alpha}
% \def\be{\beta}
% \def\ga{\gamma}
% \def\de{\delta}
% \def\ec{{{\mathsf{A}}}}
% \def\ep{\epsilon}
% \def\ka{\kappa}
% \def\lam{\lambda}
% \def\rh{\rho}
% \def\si{\sigma}
% \def\ze{\zeta}

% \def\aa{\mathsf{a}}
% \def\bb{\mathsf{b}}
% \def\cc{\mathsf{c}}
% \def\dd{\mathsf{d}}

% \def\Wh{\mathrm{W}}
% \def\Mp{M_{\text{Pl}}}
% \def\Re{\mathrm{Re}\,}
% \def\Im{\mathrm{Im}\,}

% \newcommand{\ob}[1]{\mkern 2mu \overline{\mkern -2mu #1 \mkern -2mu}\mkern 2mu}
% \newcommand{\wt}[1]{\mkern 2mu \widetilde{\mkern -2mu #1 \mkern -2mu}\mkern 2mu}
% \newcommand{\wh}[1]{\mkern 2mu \widehat{\mkern-2mu#1\mkern-2mu}\mkern 2mu}
\newcommand{\fnemail}[1]{\footnote{Email: \href{mailto:#1}{\nolinkurl{#1}}}}

% \newcommand{\zx}[1]{{\color{BrickRed}{[{ZX:} #1]}}}
% \newcommand{\jz}[1]{{\color{Blue}{[{JZ:} #1]}}}
\definecolor{CaseHeadingBG}{RGB}{235,240,248}
\newcommand{\caseheading}[1]{%
    \noindent\colorbox{CaseHeadingBG}{%
        \parbox{\dimexpr\linewidth-2\fboxsep\relax}{\bfseries #1}%
    }\par\smallskip
}

\newtheorem{theorem}{Theorem}
\newtheorem{lemma}{Lemma}

\titleformat{\section}{\normalfont\fontsize{15}{20}\bfseries\boldmath}{\thesection}{1em}{}
%%%%%%%%%%%%%%%%%%%%%%%%%%%%%%%%%%%%%%%%%%%%%%%%%%
% Use titlesec for bold text and math in headings.
\titleformat{\subsection}{\normalfont\large\bfseries\boldmath}{\thesubsection}{1em}{}
\titleformat{\subsubsection}{\normalfont\normalsize\bfseries\boldmath}{\thesubsubsection}{1em}{}
%%%%%%%%%%%%%%%%%%%%%%%%%%%%%%%%%%%%%%%%%%%%%%%%%%
\numberwithin{equation}{section}
\allowdisplaybreaks % allow page breaks inside math environments globally
%%%%%%%%%%%%%%%%%%%%%%%%%%%%%%%%%%%%%%%%%%%%%%%%%%
\usepackage{braket}
\usepackage{cancel}

\usepackage{tensor} % how to handle indices: https://tex.stackexchange.com/questions/11542/left-and-right-subscript-superscript
% left and right index: \tensor[_a]{M}{_b}

\usepackage{simpler-wick} % to use Wick contraction

\usepackage[T1]{fontenc}
\usepackage{lmodern}
%%%%%%%%%%%%%%%%%%%%%%%%%%%%%%%%%%%%%%%%%%%%%%%%%%

\newcommand{\id}{\mathbbm{1}}
\DeclareMathOperator{\Tr}{Tr}
\DeclareMathOperator{\Sym}{Sym}
\DeclareMathOperator{\Anti}{Anti}
\begin{document}
    \title{\Large\textbf{\texorpdfstring{$(A, B)$}{(A,B)} Representation of the Lorentz Group}}
    \author{Siyang Wan\fnemail{wansy25@mails.tsinghua.edu.cn} \ and ChatGPT\\[5mm] \normalsize{\emph{Department of Physics, Tsinghua University}}}
    \date{}
    \maketitle
    \tableofcontents 
    \newpage
    % Source page 1.
\section{\texorpdfstring{$(A, B)$}{(A,B)} Representation}
Start with the Lorentz Lie algebra:
\begin{equation}
    [J_i, J_j] = i \epsilon_{i j k} J_k, \qquad [K_i, K_j] = - i \epsilon_{i j k} J_k, \qquad [J_i, K_j] = i \epsilon_{i j k} K_k.
\end{equation}
Let $\boldsymbol{A} = \frac{1}{2} (\boldsymbol{J} + i \boldsymbol{K})$, $\boldsymbol{B} = \frac{1}{2} (\boldsymbol{J} - i \boldsymbol{K})$. Then
\begin{equation}
    \boldsymbol{J} = \boldsymbol{A} + \boldsymbol{B}, \qquad \boldsymbol{K} = - i (\boldsymbol{A} - \boldsymbol{B}),
\end{equation}
and we can represent $\boldsymbol{A}$ and $\boldsymbol{B}$ by
\begin{equation}
    \boldsymbol{A} = \boldsymbol{J}^{(A)} \id^B, \qquad \boldsymbol{B} = \id^A \boldsymbol{J}^{(B)}.
\end{equation}
They act on
\begin{equation}
    \ket{w} = \sum_{a, b} w_{a b} \ket{A a; B b} \in V_A \otimes V_B \equiv V_{(A, B)},
\end{equation}
explicitly,
\begin{equation}
    \begin{dcases}
        \boldsymbol{J}^{(A)} \id^B \ket{w} = \bigl(\boldsymbol{J}^{(A)}_{a a'} w_{a' b} \bigr) \ket{A a; B b}, \\
        \id^A \boldsymbol{J}^{(B)} \ket{w} = \bigl(\boldsymbol{J}^{(B)}_{b b'} w_{a b'} \bigr) \ket{A a; B b}.
    \end{dcases}
\end{equation}

    % Source pages 2--3.
\section{\texorpdfstring{$(\frac{1}{2}, 0)$ \& $(0, \frac{1}{2})$ Representations}{(1/2,0) and (0,1/2) Representations}}
\begin{framed}
    The Pauli matrices are
    \begin{equation}
        \sigma_1 = \begin{pmatrix} 0 & 1 \\
            1 & 0
        \end{pmatrix}, \qquad
        \sigma_2 = \begin{pmatrix} 0 & - i \\
            i & 0
        \end{pmatrix}, \qquad
        \sigma_3 = \begin{pmatrix} 1 & 0 \\
            0 & - 1
        \end{pmatrix}.
    \end{equation}
\end{framed}
\subsection{\texorpdfstring{$(\frac{1}{2}, 0)$ Representation}{(1/2,0) Representation}}
\begin{equation}
    V_{(\frac{1}{2}, 0)} \ni \ket{\psi} = \sum_{\alpha = \pm \frac{1}{2}} \psi_\alpha \ket{\tfrac{1}{2}, \alpha; 0, 0}
\end{equation}
and
\begin{equation}
    \begin{dcases}
        \tensor{(\boldsymbol{A}^L)}{_\alpha^\beta} = \frac{\boldsymbol{\sigma}}{2}, \\
        \tensor{(\boldsymbol{B}^L)}{_\alpha^\beta} = 0
    \end{dcases}
    \quad \Longleftrightarrow \quad
    \begin{dcases}
        \tensor{(\boldsymbol{J}^L)}{_\alpha^\beta} = \frac{\boldsymbol{\sigma}}{2}, \\
        \tensor{(\boldsymbol{K}^L)}{_\alpha^\beta} = - \frac{i \boldsymbol{\sigma}}{2},
    \end{dcases}
\end{equation}
where $L$ simply means $(\frac{1}{2}, 0)$.

\subsection{\texorpdfstring{$(0, \frac{1}{2})$ Representation}{(0,1/2) Representation}}
In the $(0, \frac{1}{2})$ representation, we have
\begin{equation}
    \begin{dcases}
        \tensor{(\boldsymbol{A}^R)}{_{\dot\alpha}^{\dot\beta}} = 0, \\
        \tensor{(\boldsymbol{B}^R)}{_{\dot\alpha}^{\dot\beta}} = - \frac{\boldsymbol{\sigma}^*}{2}
    \end{dcases}
    \quad \Longleftrightarrow \quad
    \begin{dcases}
        \tensor{(\boldsymbol{J}^R)}{_{\dot\alpha}^{\dot\beta}} = - \frac{\boldsymbol{\sigma}^*}{2}, \\
        \tensor{(\boldsymbol{K}^R)}{_{\dot\alpha}^{\dot\beta}} = - \frac{i \boldsymbol{\sigma}^*}{2}.
    \end{dcases}
\end{equation}
Define
\begin{equation}
    \epsilon_{\dot\alpha \dot\beta} = \begin{pmatrix} 0 & 1 \\
        - 1 & 0
    \end{pmatrix}, \qquad
    \epsilon^{\dot\alpha \dot\beta} = \begin{pmatrix} 0 & - 1 \\
        1 & 0
    \end{pmatrix},
\end{equation}
we have $\epsilon_{\dot\alpha \dot\beta} \epsilon^{\dot\beta \dot\gamma} = \delta_{\dot\alpha}^{\dot\gamma}$. In particular,
\begin{equation}
    \epsilon^{\dot\alpha \dot\alpha'} \tensor{(\boldsymbol{\sigma}^*)}{_{\dot\alpha'}^{\dot\beta'}} \epsilon_{\dot\beta' \dot\beta} = \tensor{(- \boldsymbol{\sigma})}{^{\dot\alpha}_{\dot\beta}},
\end{equation}
so we have
\begin{equation}
    \begin{dcases}
        \tensor{(\boldsymbol{A}^R)}{^{\dot\alpha}_{\dot\beta}} = 0, \\
        \tensor{(\boldsymbol{B}^R)}{^{\dot\alpha}_{\dot\beta}} = \frac{\boldsymbol{\sigma}}{2}
    \end{dcases}
    \quad \Longleftrightarrow \quad
    \begin{dcases}
        \tensor{(\boldsymbol{J}^R)}{^{\dot\alpha}_{\dot\beta}} = \frac{\boldsymbol{\sigma}}{2}, \\
        \tensor{(\boldsymbol{K}^R)}{^{\dot\alpha}_{\dot\beta}} = \frac{i \boldsymbol{\sigma}}{2}.
    \end{dcases}
\end{equation}
It is shown above that $\epsilon$'s help us to retain a normal $(\tfrac{1}{2})$ representation on $\boldsymbol{B}$, with
\begin{equation}
    V_{(0, \frac{1}{2})} \ni \ket{\chi} = \chi^{\dot\alpha} \ket{0, 0; \tfrac{1}{2}, \dot\alpha} = \epsilon^{\dot\alpha \dot\beta} \chi_{\dot\beta} \ket{0, 0; \tfrac{1}{2}, \dot\alpha}.
\end{equation}

    % Source pages 3--4.
\section{\texorpdfstring{Building $(\frac{1}{2}, \frac{1}{2}) = (\frac{1}{2}, 0) \otimes (0, \frac{1}{2})$}{Building (1/2,1/2) = (1/2,0) x (0,1/2)}}
\begin{equation}
    V_{(\frac{1}{2}, \frac{1}{2})} \ni \ket{v} = \tensor{v}{_\alpha^{\dot\beta}} \ket{\tfrac{1}{2}, \alpha; \tfrac{1}{2}, \dot\beta}, \quad v_{\alpha \dot{\beta}} = \epsilon_{\dot\beta \dot\beta'} \tensor{v}{_\alpha^{\dot\beta'}},
\end{equation}
and
\begin{equation}
    \tensor{(J^{\mu \nu})}{_\alpha^\beta_{\dot\gamma}^{\dot\delta}} = \tensor{(J_L^{\mu \nu})}{_\alpha^\beta} \id + \id \tensor{(J_R^{\mu \nu})}{_{\dot\gamma}^{\dot\delta}}.
\end{equation}
\subsection{Four-Vector Representation}
\begin{equation}
    V_4 = \mathbb R^{1, 3} \ni v^\mu
\end{equation}
and
\begin{equation}
    \tensor{(J^{\mu \nu})}{^\rho_\sigma} = - i \bigl(\eta^{\mu \rho} \delta^\nu_\sigma - \mu \leftrightarrow \nu \bigr) \quad \Longleftrightarrow \quad
    \begin{dcases}
        \tensor{(J_i)}{^j_k} = - i \epsilon_{i j k}, \\
        \tensor{(K_i)}{^0_j} = i \delta_{i j} = \tensor{(K_i)}{^j_0}.
    \end{dcases} \label{Four-vector representation}
\end{equation}
\subsection{The Explicit Map and the Intertwining Proof}
With the definitions
\begin{equation}
    \begin{dcases}
        \sigma^\mu_{\alpha \dot\beta} = (1, \boldsymbol{\sigma}), \\
        \bar\sigma^{\mu \dot\beta \alpha} = \epsilon^{\alpha \alpha'} \epsilon^{\dot\beta \dot\beta'} \sigma^\mu_{\alpha' \dot\beta'} = (1, - \boldsymbol{\sigma}),
    \end{dcases}
\end{equation}
the map $F:\mathbb R^{1, 3} \to V_{(\frac{1}{2}, \frac{1}{2})}$ is
\begin{equation}
    F: v^\mu \mapsto v_{\alpha \dot\beta} = v_\mu \sigma^\mu_{\alpha \dot\beta}.
\end{equation}
\begin{framed}
    The metric and trace identity are
    \begin{equation}
        \eta_{\mu \nu} = (- 1, + 1, + 1, + 1) \quad \text{and} \quad \underbrace{\Tr(\bar\sigma^\mu \sigma^\nu)}_{ \bar\sigma^{\mu \dot\beta \alpha} \sigma^\nu_{\alpha \dot\beta}} = - 2 \eta^{\mu \nu},
    \end{equation}
    So
    \begin{equation}
        v^\mu = - \frac{\bar\sigma^{\mu \dot\beta \alpha}}{2} v_{\alpha \dot\beta} = - \frac{1}{2} \Tr(\bar\sigma^\mu v).
    \end{equation}
\end{framed}
It can be shown that $v^\mu$ and $v_{\alpha \dot\beta}$ transform equivalently under the Lorentz group, see Eq. \eqref{Four-vector representation of the Lorentz transformation}.

    % Source pages 5--6.
\section{\texorpdfstring{$(\frac{1}{2}, \frac{1}{2}) \otimes (\frac{1}{2}, \frac{1}{2})$: Rank-Two Tensors}{(1/2,1/2) x (1/2,1/2): Rank-Two Tensors}}
\begin{equation}
    V_{(\frac{1}{2}, \frac{1}{2}) \otimes (\frac{1}{2}, \frac{1}{2})} \ni \ket{T} = \tensor{T}{_\alpha_\gamma^{\dot\beta}^{\dot\delta}} \ket{\tfrac{1}{2}, \alpha; \tfrac{1}{2}, \dot\beta} \ket{\tfrac{1}{2}, \gamma; \tfrac{1}{2}, \dot\delta}
\end{equation}
and
\begin{equation}
    T_{\alpha \gamma \dot\beta \dot\delta} = T_{\mu \nu} \sigma^\mu_{\alpha \dot\beta} \sigma^\nu_{\gamma \dot\delta}.
\end{equation}
\subsection{\texorpdfstring{$(\frac{1}{2}, \frac{1}{2}) \otimes (\frac{1}{2}, \frac{1}{2}) = (1, 1) \oplus (1, 0) \oplus (0, 1) \oplus (0, 0)$}{(1/2,1/2) x (1/2,1/2) = (1,1) + (1,0) + (0,1) + (0,0)}}
First, $(\frac{1}{2}, 0) \otimes (\frac{1}{2}, 0) = (1, 0) \oplus (0, 0)$ and
\begin{align}
    \ket{1, 1}_L & = \ket{\tfrac{1}{2}, \tfrac{1}{2}; \tfrac{1}{2}, \tfrac{1}{2}}_L, \notag \\
    \ket{1, 0}_L & = \frac{1}{\sqrt{2}} \bigl( \ket{\tfrac{1}{2}, \tfrac{1}{2}; \tfrac{1}{2}, - \tfrac{1}{2}}_L + \ket{\tfrac{1}{2}, - \tfrac{1}{2}; \tfrac{1}{2}, \tfrac{1}{2}}_L \bigr), \notag \\
    \ket{1, - 1}_L & = \ket{\tfrac{1}{2}, - \tfrac{1}{2}; \tfrac{1}{2}, - \tfrac{1}{2}}_L, \notag \\[2mm]
    \ket{0, 0}_L & = \frac{1}{\sqrt{2}} \bigl( \ket{\tfrac{1}{2}, \tfrac{1}{2}; \tfrac{1}{2}, - \tfrac{1}{2}}_L - \ket{\tfrac{1}{2}, - \tfrac{1}{2}; \tfrac{1}{2}, \tfrac{1}{2}}_L \bigr).
\end{align}
Similarly, $(0, \frac{1}{2}) \otimes (0, \frac{1}{2}) = (0, 1) \oplus (0, 0)$.

The CG coefficients are
\begin{equation}
    \braket{\tfrac{1}{2}, \alpha; \tfrac{1}{2}, \gamma | 0, 0}_L = \frac{\epsilon_{\alpha \gamma}}{\sqrt{2}}, \qquad \braket{\tfrac{1}{2}, \dot\beta; \tfrac{1}{2}, \dot\delta | 0, 0}_R = \frac{\epsilon_{\dot\beta \dot\delta}}{\sqrt{2}}.
\end{equation}
With the CG coefficients, we can write
\begin{equation}
    \ket{T} = \sum_{j = 0}^1 \sum_{k = 0}^1 T_{j m, k n} \ket{j m}_L \otimes \ket{k n}_R,
\end{equation}
and
\begin{equation}
    T_{\alpha \gamma \dot\beta \dot\delta} = \sum_{j m k n} T_{j m, k n} \braket{\tfrac{1}{2}, \alpha; \tfrac{1}{2}, \gamma | j m}_L \braket{\tfrac{1}{2}, \dot\beta; \tfrac{1}{2}, \dot\delta | k n}_R. \label{From (A,B) to spinor indices}
\end{equation}
\subsection{\texorpdfstring{$(0, 0)$}{(0,0)}}
The $j = k = m = n = 0$ part of the tensor is
\begin{align}
    & T_{\alpha \gamma \dot\beta \dot\delta} \supset T_{0 0 0 0} \frac{\epsilon_{\alpha \gamma} \epsilon_{\dot\beta \dot\delta}}{2} \qquad (\text{let }T_0 \equiv T_{0 0 0 0}) \\
    \Rightarrow \quad & T_{\mu \nu} \supset \frac{T_0}{8} \Tr(\bar\sigma_\mu \sigma_\nu) = - \frac{T_0}{4} \eta_{\mu \nu}.
\end{align}
\subsection{\texorpdfstring{$(1, 0)$ \& $(0, 1)$}{(1,0) and (0,1)}}
First, we consider $j = 1$, $k = 0$:
\begin{equation}
    T_{\alpha \gamma \dot\beta \dot\delta} \supset T_{1 m 0 0} \braket{\tfrac{1}{2}, \alpha; \tfrac{1}{2}, \gamma | 1 m} \frac{\epsilon_{\dot\beta \dot\delta}}{\sqrt{2}} \qquad (\text{let }T_{1 m 0 0} \equiv T_m^L).
\end{equation}
\begin{framed}
    A crucial observation here is that
    \begin{align}
        & \braket{\tfrac{1}{2}, \alpha; \tfrac{1}{2}, \gamma | 1 m} = \Sym_{\alpha \gamma} \braket{\cdots}, \\
        & \Sym_{\alpha \gamma} T_{\alpha \gamma \cdots} = \frac{T_{\alpha \gamma \cdots} + T_{\gamma \alpha \cdots}}{2}
    \end{align}
    and
    \begin{equation}
        \braket{\tfrac{1}{2}, \alpha; \tfrac{1}{2}, \gamma | 0 0} = \epsilon_{\alpha \gamma} \left( - \frac{\epsilon^{\alpha' \gamma'}}{2} \braket{\tfrac{1}{2}, \alpha'; \tfrac{1}{2}, \gamma' | 0 0} \right).
    \end{equation}
\end{framed}
These operations project onto the $(1, 1)$, $(1, 0)$, $(0, 1)$, and $(0, 0)$ representations.

Explicitly,
\begin{align}
    \Sym_{\alpha \gamma} \Sym_{\dot\beta \dot\delta} T_{\alpha \gamma \dot\beta \dot\delta} & \in (1, 1), \notag \\
    \Sym_{\alpha \gamma} \Anti_{\dot\beta \dot\delta} T & \in (1, 0), \notag \\
    \Anti_{\alpha \gamma} \Sym_{\dot\beta \dot\delta} T & \in (0, 1), \notag \\
    \Anti_{\alpha \gamma} \Anti_{\dot\beta \dot\delta} T & \in (0, 0).
\end{align}

    % Source pages 7--10.
\section{Projection Operators}
Since $T_{\alpha \gamma \dot\beta \dot\delta} = T_{\mu \nu} \sigma^\mu_{\alpha \dot\beta} \sigma^\nu_{\gamma \dot\delta}$, we can write down the projected tensor using $\mu \nu$ indices.
\begin{framed}
    The inner product is
    \begin{equation}
        u \cdot v = - \frac{1}{2} \epsilon^{\dot\beta \dot\beta'} \epsilon^{\alpha \alpha'} u_{\alpha' \dot\beta'} v_{\alpha \dot\beta}.
    \end{equation}
\end{framed}
The projection operators are
\begin{equation}
    \tensor{{P^{(1)}}}{_a^{a'}_b^{b'}} = \frac{1}{2} (\delta_a^{a'} \delta_b^{b'} + \delta_b^{a'} \delta_a^{b'}), \qquad \tensor{{P^{(0)}}}{_a^{a'}_b^{b'}} = - \frac{1}{2} \epsilon_{a b} \epsilon^{a' b'}
\end{equation}
and
\begin{align}
    \tensor{{P^{(1, 1)}}}{_\alpha^{\alpha'}_\gamma^{\gamma'}_{\dot\beta}^{\dot\beta'}_{\dot\delta}^{\dot\delta'}} & = \tensor{{P^{(1)}}}{_\alpha^{\alpha'}_\gamma^{\gamma'}} \tensor{{P^{(1)}}}{_{\dot\beta}^{\dot\beta'}_{\dot\delta}^{\dot\delta'}}, \notag \\
    P^{(1, 0)} & = P^{(1)} \otimes P^{(0)}, \notag \\
    & \cdots
\end{align}
\begin{framed}
    \begin{align}
        \tensor{P}{_\mu^{\mu'}_\nu^{\nu'}} = {} & \left(- \frac{\bar\sigma_\mu^{\dot\beta \alpha}}{2} \right) \left(- \frac{\bar\sigma^{\mu' \dot\beta'' \alpha''}}{2} \underbrace{\epsilon_{\dot\beta'' \dot\beta'} \epsilon_{\alpha'' \alpha'}}_{\text{contribute } (- 2)} \times (- 2) \right) \notag \\
        & \times \left(- \frac{\bar\sigma_\nu^{\dot\delta \gamma}}{2} \right) \left(- \frac{\sigma^{\nu'}_{\gamma' \dot\delta'}}{2} \times (- 2) \right) \tensor{P}{_\alpha^{\alpha'}_\gamma^{\gamma'}_{\dot\beta}^{\dot\beta'}_{\dot\delta}^{\dot\delta'}}.
    \end{align}
\end{framed}
\begin{framed}
    More naturally, we should raise and lower indices with
    \begin{equation}
        \frac{i}{\sqrt{2}} \epsilon^{\alpha \beta} \quad \& \quad \frac{\sqrt{2}}{i} \epsilon_{\alpha \beta}.
    \end{equation}
\end{framed}
The indices of the preceding $P^{\alpha' \gamma' \dot\beta' \dot\delta'}$ are \emph{not} raised by $\epsilon^{\alpha \beta}$ but by $\frac{\sqrt{2}}{i} \epsilon^{\alpha \beta}$.
% Source page 8.
Consider
\begin{align}
    L & = (P^{(1)} - P^{(0)}) \otimes (P^{(1)} + P^{(0)}) \notag \\
    & = P^{(1, 1)} + P^{(1, 0)} - P^{(0, 1)} - P^{(0, 0)}
\end{align}
and
\begin{align}
    R & = (P^{(1)} + P^{(0)}) \otimes (P^{(1)} - P^{(0)}) \notag \\
    & = P^{(1, 1)} - P^{(1, 0)} + P^{(0, 1)} - P^{(0, 0)},
\end{align}
Conversely,
\begin{equation}
    P^{(j, k)} = \frac{1}{4} (I + s_j L) (I + s_k R) \quad \text{with} \quad
    \begin{dcases}
        s_1 = + 1, \\
        s_0 = - 1.
    \end{dcases}
\end{equation}
\begin{framed}
    \begin{align}
        L R & = (P^{(1)} - P^{(0)}) \otimes (P^{(1)} - P^{(0)}) \notag \\
        & = P^{(1, 1)} - P^{(1, 0)} - P^{(0, 1)} + P^{(0, 0)}
    \end{align}
    since
    \begin{equation}
        P^{(1)} \cdot P^{(1)} = P^{(1)}, \qquad P^{(0)} \cdot P^{(0)} = P^{(0)}
    \end{equation}
    and
    \begin{equation}
        P^{(1)} \cdot P^{(0)} = P^{(0)} \cdot P^{(1)} = 0.
    \end{equation}
\end{framed}
The benefit of considering $L$ and $R$ is that
\begin{equation}
    \begin{dcases}
        \tensor{(P^{(1)} + P^{(0)})}{_a^{a'}_b^{b'}} = \delta_a^{a'} \delta_b^{b'}, \\
        \tensor{(P^{(1)} - P^{(0)})}{_a^{a'}_b^{b'}} = \delta_a^{b'} \delta_b^{a'}.
    \end{dcases}
\end{equation}
% Source page 9. The displayed L/R signs below follow the manuscript.
This makes $L$ and $R$ easy to project to tensor indices,
\begin{align}
    \tensor{L}{_\mu^{\mu'}_\nu^{\nu'}} & = \frac{1}{4} \bar\sigma_\mu^{\dot\beta \alpha} \sigma^{\mu'}_{\alpha' \dot\beta'} \bar\sigma_\nu^{\dot\delta \gamma} \sigma^{\nu'}_{\gamma' \dot\delta'} \delta_\alpha^{\gamma'} \delta_\gamma^{\alpha'} \delta_{\dot\beta}^{\dot\beta'} \delta_{\dot\delta}^{\dot\delta'} \notag \\
    & = \frac{1}{4} \Tr(\sigma^{\nu'} \bar\sigma_\nu \sigma^{\mu'} \bar\sigma_\mu) \notag \\
    & = \frac{1}{2} \bigl(\delta_\mu^{\mu'} \delta_\nu^{\nu'} + \mu' \leftrightarrow \nu' - \eta_{\mu \nu} \eta^{\mu' \nu'} + i \tensor{\epsilon}{_{\mu \nu}^{\mu' \nu'}} \bigr). \\
    \tensor{R}{_\mu^{\mu'}_\nu^{\nu'}} & = \frac{1}{4} \bar\sigma_\mu^{\dot\beta \alpha} \sigma^{\mu'}_{\alpha' \dot\beta'} \bar\sigma_\nu^{\dot\delta \gamma} \sigma^{\nu'}_{\gamma' \dot\delta'} \delta_\alpha^{\alpha'} \delta_\gamma^{\gamma'} \delta_{\dot\beta}^{\dot\delta'} \delta_{\dot\delta}^{\dot\beta'} \notag \\
    & = \frac{1}{4} \Tr(\bar\sigma_\mu \sigma^{\mu'} \bar\sigma_\nu \sigma^{\nu'}) \notag \\
    & = \tensor{L}{_\mu^{\nu'}_\nu^{\mu'}} \notag \\
    & = \frac{1}{2} \bigl(\delta_\mu^{\mu'} \delta_\nu^{\nu'} + \mu' \leftrightarrow \nu' - \eta_{\mu \nu} \eta^{\mu' \nu'} - i \tensor{\epsilon}{_{\mu \nu}^{\mu' \nu'}} \bigr).
\end{align}
A more compact notation is
\begin{equation}
    I = \delta_\mu^{\mu'} \delta_\nu^{\nu'}, \qquad X = \delta_\mu^{\nu'} \delta_\nu^{\mu'}, \qquad G = \eta_{\mu \nu} \eta^{\mu' \nu'}
\end{equation}
then
\begin{equation}
    L = \frac{1}{2} (I + X - G + i \epsilon), \qquad R = \frac{1}{2} (I + X - G - i \epsilon)
\end{equation}
and
\begin{equation}
    P^{(j, k)} = \frac{1}{4} (I + s_j L + s_k R + s_j s_k X).
\end{equation}
% Source page 10.
Explicitly,
\begin{align}
    P^{(1, 1)} & = \frac{1}{2} (I + X) - \frac{1}{4} G, \notag \\
    P^{(1, 0)} & = \frac{1}{4} (I - X + i \epsilon), \notag \\
    P^{(0, 1)} & = \frac{1}{4} (I - X - i \epsilon), \notag \\
    P^{(0, 0)} & = \frac{1}{4} G.
\end{align}
For a tensor $T_{\mu \nu}$, let $T = \tensor{T}{^\mu_\mu}$ and
\begin{equation}
    \begin{dcases}
        F_{\mu \nu} = T_{[\mu \nu]}, \\
        (* F)_{\mu \nu} = \frac{1}{2} \tensor{\epsilon}{_{\mu \nu}^{\rho \sigma}} F_{\rho \sigma},
    \end{dcases}
\end{equation}
then
\begin{framed}
    \begin{align}
        (P^{(1, 1)} T)_{\mu \nu} & = T_{(\mu \nu)} - \frac{1}{4} \eta_{\mu \nu} T, \notag \\
        (P^{(1, 0)} T)_{\mu \nu} & = \frac{1}{2} \bigl(F_{\mu \nu} + i (* F)_{\mu \nu} \bigr), \notag \\
        (P^{(0, 1)} T)_{\mu \nu} & = \frac{1}{2} \bigl(F_{\mu \nu} - i (* F)_{\mu \nu} \bigr), \notag \\
        (P^{(0, 0)} T)_{\mu \nu} & = \frac{1}{4} \eta_{\mu \nu} T.
    \end{align}
\end{framed}

    \section{From \texorpdfstring{$(A, B)$}{(A,B)} to Spin}
We restrict to the massive little group
\begin{equation}
    (A, B) \Big|_{\mathrm{SO}(3)} = \bigoplus_{j = |A - B|}^{A + B} (j),
\end{equation}
choose a particle spin $j$, obtain the CG coefficients
\begin{equation}
    u_{a b}(k, \sigma) = \braket{A a; B b | j \sigma},
\end{equation}
and boost them to
\begin{equation}
    u_{\mu \nu}(p, \sigma)
\end{equation}
to build the spin-$j$ covariant field. See \textit{1PSs \& Field Operators} for more detail.

As a more concrete example, let us consider
\begin{equation}
    (1, 1) \Big|_{\mathrm{SO}(3)} = (0) \oplus (1) \oplus (2).
\end{equation}
First we need to project $a, b$ to the spinor indices by \eqref{From (A,B) to spinor indices}, then project the spinor indices to tensor indices, explicitly
\begin{align}
    & \begin{dcases}
        \tensor{u}{_\alpha_\gamma^{\dot\beta}^{\dot\delta}}(k; j, \sigma) = u_{a b}(k; j, \sigma) \braket{\tfrac{1}{2} \alpha; \tfrac{1}{2} \gamma | 1 a} \braket{\tfrac{1}{2} \dot{\beta}; \tfrac{1}{2} \dot{\delta} | 1 b}, \\
        u_{\mu \nu}(k; j, \sigma) = \frac{1}{4} \bar{\sigma}^{\dot{\beta} \alpha}_\mu \bar{\sigma}^{\dot{\delta} \gamma}_\nu u_{\alpha \gamma \dot{\beta} \dot{\delta}}(k; j, \sigma)
    \end{dcases} \\
    \Rightarrow \quad & P^{a b}_{\mu \nu} = \frac{1}{4} \bar{\sigma}^{\dot{\beta} \alpha}_\mu \bar{\sigma}^{\dot{\delta} \gamma}_\nu \epsilon_{\dot{\beta} \dot{\beta}'} \epsilon_{\dot{\delta} \dot{\delta}'} \braket{\tfrac{1}{2} \alpha; \tfrac{1}{2} \gamma | 1 a} \braket{\tfrac{1}{2} \dot{\beta}'; \tfrac{1}{2} \dot{\delta}' | 1 b}.
\end{align}
We obtain the following rest-frame, $k = (m, 0, 0, 0)$, coefficient tensors. For $j = 0$,
\begin{equation}
    u_{\mu \nu}(k; 0, 0) = \frac{1}{4 \sqrt{3}} \begin{pmatrix}
        3 & 0 & 0 & 0 \\
        0 & 1 & 0 & 0 \\
        0 & 0 & 1 & 0 \\
        0 & 0 & 0 & 1
    \end{pmatrix}.
\end{equation}
For $j = 1$,
\begin{align}
    u_{\mu \nu}(k; 1, 1) &= \frac{1}{4} \begin{pmatrix}
        0 & - 1 & - i & 0 \\
        - 1 & 0 & 0 & 0 \\
        - i & 0 & 0 & 0 \\
        0 & 0 & 0 & 0
    \end{pmatrix}, \quad u_{\mu \nu}(k; 1, 0) = \frac{1}{2 \sqrt{2}} \begin{pmatrix}
        0 & 0 & 0 & 1 \\
        0 & 0 & 0 & 0 \\
        0 & 0 & 0 & 0 \\
        1 & 0 & 0 & 0
    \end{pmatrix}, \\
    u_{\mu \nu}(k; 1, - 1) &= \frac{1}{4} \begin{pmatrix}
        0 & 1 & - i & 0 \\
        1 & 0 & 0 & 0 \\
        - i & 0 & 0 & 0 \\
        0 & 0 & 0 & 0
    \end{pmatrix}.
\end{align}
For $j = 2$,
\begin{align}
    u_{\mu \nu}(k; 2, 2) &= \frac{1}{4} \begin{pmatrix}
        0 & 0 & 0 & 0 \\
        0 & 1 & i & 0 \\
        0 & i & - 1 & 0 \\
        0 & 0 & 0 & 0
    \end{pmatrix}, \quad u_{\mu \nu}(k; 2, 1) = \frac{1}{4} \begin{pmatrix}
        0 & 0 & 0 & 0 \\
        0 & 0 & 0 & - 1 \\
        0 & 0 & 0 & - i \\
        0 & - 1 & - i & 0
    \end{pmatrix}, \\
    u_{\mu \nu}(k; 2, 0) &= \frac{1}{2 \sqrt{6}} \begin{pmatrix}
        0 & 0 & 0 & 0 \\
        0 & - 1 & 0 & 0 \\
        0 & 0 & - 1 & 0 \\
        0 & 0 & 0 & 2
    \end{pmatrix}, \quad u_{\mu \nu}(k; 2, - 1) = \frac{1}{4} \begin{pmatrix}
        0 & 0 & 0 & 0 \\
        0 & 0 & 0 & 1 \\
        0 & 0 & 0 & - i \\
        0 & 1 & - i & 0
    \end{pmatrix}, \\
    u_{\mu \nu}(k; 2, - 2) & = \frac{1}{4} \begin{pmatrix}
        0 & 0 & 0 & 0 \\
        0 & 1 & - i & 0 \\
        0 & - i & - 1 & 0 \\
        0 & 0 & 0 & 0
    \end{pmatrix}.
\end{align}
All nine tensors are symmetric and traceless.
Under the massive little group $\mathrm{SO}(3)$, their components are organized as
\begin{align}
    j = 0 &: \quad \boxed{u_{0 0}}, \quad u_{i j} = \frac{1}{3} \delta_{i j} u_{0 0}, \notag \\
    j = 1 &: \quad \boxed{u_{0 i}} = u_{i 0}, \quad u_{0 0} = u_{i j} = 0, \notag \\
    j = 2 &: \quad \boxed{u_{i j}} = u_{j i}, \quad \delta^{i j} u_{i j} = 0, \quad u_{0 0} = u_{0 i} = 0. \label{From (1,1) to spins}
\end{align}

In particular, consider
\begin{equation}
    O_{\mu \nu} = \Big( \partial_\mu \partial_\nu - \frac{1}{4} \eta_{\mu \nu} \partial^2 \Big) \sigma,
\end{equation}
we have
\begin{equation}
    u_{\mu \nu}(k) \propto \braket{0 | O_{\mu \nu}(0) | k} = - \Big( k_\mu k_\nu - \frac{1}{4} \eta_{\mu \nu} k^2 \Big) \braket{0 | \sigma(0) | k}.
\end{equation}
Comparing to \eqref{From (1,1) to spins} we know $O_{\mu \nu}$ creates spin-0 particles.

    \section{\texorpdfstring{$(\frac{1}{2}, \frac{1}{2})$}{(1/2,1/2)} and Clifford Algebra}
The complexified Clifford algebra $\mathrm{Cl}_{1,3}(\mathbb R)_{\mathbb C}$ has the grades
\begin{align}
    \text{grade 0:} \quad & I && \text{(scalar)}, \notag \\
    \text{grade 1:} \quad & \gamma^\mu && \text{(vector)}, \notag \\
    \text{grade 2:} \quad & J^{\mu \nu} = - \frac{i}{4} [\gamma^\mu, \gamma^\nu] && \text{(bivector/tensor)}, \notag \\
    \text{grade 3:} \quad & \gamma^\mu \gamma_5 && \text{(pseudovector)}, \notag \\
    \text{grade 4:} \quad & \gamma_5 = - i \gamma^0 \gamma^1 \gamma^2 \gamma^3 && \text{(pseudoscalar)},
\end{align}
with
\begin{equation}
    \{\gamma^\mu, \gamma^\nu\} = 2 \eta^{\mu \nu}, \qquad \eta^{\mu \nu} = (- 1, + 1, + 1, + 1).
\end{equation}
\begin{framed}
    It can be shown that
    \begin{equation}
        \begin{dcases}
            [J^{\mu \nu}, J^{\rho \sigma}] = i (\eta^{\mu \rho} J^{\nu \sigma} - \cdots), \\
            [J^{\mu \nu}, \gamma^\rho] = i (\eta^{\mu \rho} \gamma^\nu - \eta^{\nu \rho} \gamma^\mu).
        \end{dcases}
    \end{equation}
\end{framed}
\begin{framed}
    Although not immediately related to these notes, the parity transformation is related to $\beta = i \gamma^0$, with
    \begin{align}
        \beta \gamma^i \beta^{- 1} & = - \gamma^i, \qquad \beta \gamma^0 \beta^{- 1} = \gamma^0, \notag \\
        \beta J^{i j} \beta^{- 1} & = J^{i j}, \qquad \beta J^{i 0} \beta^{- 1} = - J^{i 0}, \qquad \beta \gamma_5 \beta^{- 1} = - \gamma_5.
    \end{align}
\end{framed}

% Source page 12. The corrected handwritten heading uses L direct-sum R.
\subsection{\texorpdfstring{$L \oplus R$ and $\mathrm{Cl}_{1,3}$ from $\gamma_5$}{L + R and Cl(1,3) from gamma5}}
Observe that
\begin{equation}
    \begin{dcases}
        [\gamma_5, J^{\mu \nu}] = 0, \qquad \text{simultaneous diagonalization for a specific } \mu \nu, \\
        (\gamma_5)^2 = I, \qquad \text{eigenvalues are } \pm 1.
    \end{dcases}
\end{equation}
\begin{framed}
    Not directly relevant here, but
    \begin{equation}
        \{\gamma_5, \gamma^\mu\} = 0.
    \end{equation}
\end{framed}
We can construct $\ket{\psi_L} \oplus \ket{\psi_R}$, or
\begin{equation}
    \begin{pmatrix}
        \psi_L \\
        \psi_R
    \end{pmatrix},
\end{equation}
and
\begin{equation}
    \gamma_5 = \begin{pmatrix}
        I_L & 0 \\
        0 & - I_R
    \end{pmatrix}.
\end{equation}
\subsection{\texorpdfstring{Dirac $(\frac{1}{2}, 0) \oplus (0, \frac{1}{2})$ Representation}{Dirac (1/2,0) + (0,1/2) Representation}}
\begin{align}
    \gamma^\mu & = - i \begin{pmatrix}
        \tensor{0}{_\alpha^\beta} & (\sigma^\mu)_{\alpha \dot\beta} \\
        \bar\sigma^{\mu \dot\alpha \beta} & \tensor{0}{^{\dot\alpha}_{\dot\beta}}
    \end{pmatrix}, \\
    J^{\mu \nu} & = \frac{i}{4} \begin{pmatrix}
        \tensor{(\sigma^{\mu \nu})}{_\alpha^\beta} & 0_{\alpha \dot\beta} \\
        0^{\dot\alpha \beta} & \tensor{(\bar\sigma^{\mu \nu})}{^{\dot\alpha}_{\dot\beta}}
    \end{pmatrix},
\end{align}
where
\begin{equation}
    \begin{dcases}
        \sigma^{\mu \nu} = \sigma^\mu \bar\sigma^\nu - \sigma^\nu \bar\sigma^\mu, \\
        \bar\sigma^{\mu \nu} = \bar\sigma^\mu \sigma^\nu - \bar\sigma^\nu \sigma^\mu.
    \end{dcases}
\end{equation}
\begin{framed}
    \begin{equation}
        \gamma^\mu: L \to R, \quad R \to L.
    \end{equation}
\end{framed}

% Source page 13. Crossed-out content is omitted.
Since $\{\gamma^\mu, \gamma^\nu\} = 2 \eta^{\mu \nu}$, we have
\begin{equation}
    \begin{dcases}
        \sigma^\mu \bar\sigma^\nu + \sigma^\nu \bar\sigma^\mu = - 2 \eta^{\mu \nu} \tensor{\id}{_\alpha^\beta}, \\
        \bar\sigma^\mu \sigma^\nu + \bar\sigma^\nu \sigma^\mu = - 2 \eta^{\mu \nu} \tensor{\id}{^{\dot\alpha}_{\dot\beta}},
    \end{dcases}
\end{equation}
or
\begin{equation}
    \sigma^\mu_{\alpha \dot\gamma} \sigma^\nu_{\beta \dot\delta} \epsilon^{\dot\gamma \dot\delta} + \mu \leftrightarrow \nu = - 2 \eta^{\mu \nu} \epsilon_{\beta \alpha}.
\end{equation}
\subsection{\texorpdfstring{$(\frac{1}{2}, \frac{1}{2})$ Representation}{(1/2,1/2) Representation}}
\begin{align}
    & \tensor{\left(\exp\left(i \frac{\omega_{\mu \nu}}{2} \frac{i}{4} \sigma^{\mu \nu}\right)\right)}{_\alpha^\beta} \epsilon_{\dot\alpha \dot\alpha'} \tensor{\left(\exp\left(i \frac{\omega_{\mu \nu}}{2} \frac{i}{4} \bar\sigma^{\mu \nu}\right)\right)}{^{\dot\alpha'}_{\dot\beta'}} \epsilon^{\dot\beta' \dot\beta} \sigma^\mu_{\beta \dot\beta} \notag \\
    & = \tensor{\left(\exp\left(- \frac{\omega_{\rho \sigma}}{2} i J^{\rho \sigma}\right)\right)}{^\mu_\nu} \sigma^\nu_{\alpha \dot\alpha}. \label{Four-vector representation of the Lorentz transformation}
\end{align}
The minus sign ensures that $v_\mu \sigma^\mu_{\alpha \dot\alpha}$ and $v_\mu$ transform correctly.

\begin{framed}
    Let $\omega_{\mu \nu}$ to be infinitesimal,
    \begin{align}
        \text{LHS} &= \Big( 1 + \frac{- 1}{8} \omega_{\mu \nu} \sigma^{\mu \nu} \Big) \sigma^\mu \Big( 1 - \frac{- 1}{8} \omega_{\mu \nu} \bar{\sigma}^{\mu \nu} \Big) \notag \\
        &= \sigma^\mu - \frac{1}{8} \omega_{\rho \sigma} (\sigma^{\rho \sigma} \sigma^\mu - \sigma^\mu \bar{\sigma}^{\rho \sigma}).
    \end{align}
    The Clifford algebra tells us
    \begin{align}
        & [J^{\mu \nu}, \gamma^\rho] = i (\eta^{\mu \rho} \gamma^\nu - \eta^{\nu \rho} \gamma^\mu) \notag \\
        \Rightarrow \quad & \frac{1}{4} \begin{pmatrix}
            0 & \sigma^{\mu \nu} \sigma^\rho - \sigma^\rho \bar{\sigma}^{\mu \nu} \\
            \bar{\sigma}^{\mu \nu} \bar{\sigma}^\rho - \bar{\sigma}^\rho \sigma^{\mu \nu} & 0
        \end{pmatrix} = \begin{pmatrix}
            0 & \eta^{\mu \rho} \sigma^\nu - \eta^{\nu \rho} \sigma^\mu \\
            \eta^{\mu \rho} \bar{\sigma}^\nu - \eta^{\nu \rho} \bar{\sigma}^\mu & 0
        \end{pmatrix} \notag \\
        \Rightarrow \quad & - \frac{1}{8} \omega_{\rho \sigma} (\sigma^{\rho \sigma} \sigma^\mu - \sigma^\mu \bar{\sigma}^{\rho \sigma}) = - \frac{1}{2} \omega_{\rho \sigma} (\eta^{\rho \mu} \sigma^\sigma - \eta^{\sigma \mu} \sigma^\rho).
    \end{align}
    We have
    \begin{equation}
        \text{LHS} = \sigma^\mu - \frac{1}{2} (\tensor{\omega}{^\mu_\nu} - \tensor{\omega}{_\nu^\mu}) \sigma^\nu = \text{RHS},
    \end{equation}
    where we used \eqref{Four-vector representation},
    \begin{equation}
        - \frac{\omega_{\rho \sigma}}{2} i \tensor{(J^{\rho \sigma})}{^\mu_\nu} = - \frac{1}{2} (\tensor{\omega}{^\mu_\nu} - \tensor{\omega}{_\nu^\mu}).
    \end{equation}
\end{framed}


\end{document}
